% ======================================================
% Appendix: controlled cohort evidence
% ======================================================

\section{Controlled Cohort Evidence}
\label{app:controlled}

This appendix contains the evidence tables for the controlled R1--R7
cohort. It replaces the earlier long source note that centered on the old
Cake2 trace. The old Cake2 artifact is still preserved in the viewer
artifact, but it is not used in any quantitative table here.

\subsection{Source-Anchor Convention}
\label{app:anchor-convention}

All source anchors refer to upstream OpenClaw v2026.7.1-2 at commit
\seqsplit{0790d9f593ad30c940ed93b5872a8cf6d6f3cf8c}. When a runtime
event is emitted by TraceClaw instrumentation, the cited source decision
site is still the upstream branch or function being observed, not the
patched instrumentation line. This convention keeps source claims stable
even when the instrumentation patch adds event-emission code nearby.

The instrumentation patch hash used for the controlled R6/R7 collection
was
\seqsplit{6736a6be821ae927ebb112328d9552386684d5cbd8090975f52cc8666d8ef11f}.
The patch is treated as a supplementary artifact and should be sanitized
before submission.

\subsection{Run Provenance}
\label{app:provenance}

R1--R5 are saved trace commits from the repository history. R6 and R7
were collected in the temporary worktree
\texttt{traceclaw-r6-r7-worktree-20260930} on branch
\texttt{collect-r6-r7}. They were committed as trace-data commits only:
\texttt{data/cases/latest-live.js}, \texttt{data/cases/recent-runs.js},
and \texttt{index.html}. No collector, instrumentation, schema, UI, or
configuration code changed between 7d742df and the R7 commit.

The traces themselves do not record \texttt{openclawCommit},
\texttt{traceclawCommit}, or \texttt{configFingerprint}. The paper
therefore uses the wording ``OpenClaw v2026.7.1-2 (0790d9f) with a
TraceClaw instrumentation patch applied'' rather than claiming that each
trace embeds complete environment metadata.

\begin{table}[H]
\centering
\caption{Controlled R1--R7 cohort.}
\label{tab:app-controlled-cohort}
{\scriptsize
\begin{tabularx}{\linewidth}{@{}l X l X c c@{}}
\toprule
Run & Prompt class & Run ID & Tool summary & Gateway & AR \\
\midrule
R1 & no-tool & \texttt{8d07c88e} & none & 19/19, 19/19 & observed \\
R2 & read-only retrieval & \texttt{fba117de} & \texttt{web\_search} \( \times 6 \) & 19/19, 19/19 & observed \\
R3 & external side effect via host app & \texttt{5eb26fda} & \texttt{bash} \( \times 5 \) & 19/19, 19/19 & observed \\
R4 & local file write & \texttt{37bf6523} & \texttt{bash} \( \times 3 \) + \texttt{apply\_patch} & 19/19, 19/19 & observed \\
R5 & no-tool & \texttt{42b573dc} & none & 19/19, 19/19 & observed \\
R6 & no-tool & \texttt{3eb51151} & none & 19/19, 19/19 & observed \\
R7 & no-tool & \texttt{02225265} & none & 19/19, 19/19 & observed \\
\bottomrule
\end{tabularx}

}
\end{table}

\section{R1 Gateway Ledger}
\label{app:r1-ledger}

R1 is the primary controlled case for the paper. It uses shared-token
authentication, reaches G1 \texttt{allow}, finishes G2 with
\texttt{operator.write}, and authorizes \texttt{chat.send} through the
ordinary write check. The old Cake2 device-token/operator-admin path is
not imported into this case.

\begin{longtable}{@{}p{0.06\linewidth}p{0.14\linewidth}p{0.15\linewidth}p{0.16\linewidth}p{0.41\linewidth}@{}}
\caption{R1 per-stage Gateway ledger. Values are redacted or shortened where they are identifiers; full raw values remain in the saved trace artifact.}\label{tab:app-r1-ledger}\\
\toprule
Stage & Branch & Outcome evidence & Rationale visibility & R1 value / note \\
\midrule
\endfirsthead
\toprule
Stage & Branch & Outcome evidence & Rationale visibility & R1 value / note \\
\midrule
\endhead
G0 & authorized & Observed & partial & scope = gateway-runtime; authMode = token; authMethod = token; authOk = True; sharedAuthOk = True; sharedAuthProvided = True; hasDeviceIdentity = False; result = authorized \\
G1 & allow & Observed & partial & scope = gateway-runtime; authMode = token; authMethod = token; sharedAuthProvided = True; sharedAuthOk = True; result = allow \\
G2 & pass & Observed & partial & scope = gateway-runtime; authMode = token; authMethod = token; role = operator; scopes = [operator.write]; result = pass \\
G3 & allow & Observed & partial & result = allow \\
G4 & pass & Observed & partial & result = pass; hasPrivilegedFields = False; hasExplicitOrigin = False \\
G5 & unchanged & Observed & partial & result = unchanged; messageLength = 19; messageChangedBySanitization = False \\
G6 & resolved & Observed & partial & requestedAgentId = none; result = resolved \\
G7 & resolved & Observed & partial & canonicalSessionKey = agent:main:dashboard:trace-...; result = resolved \\
G8 & pass & Observed & partial & result = pass \\
G9 & resolved & Observed & partial & agentId = main; result = resolved \\
G10 & allow & Observed & partial & sendPolicy = allow; result = allow \\
G11 & pass & Observed & partial & dedupeDecision = new\_dispatch; result = pass \\
G12 & admitted & Observed & partial & admissionDecision = admitted; result = admitted \\
G13 & constructed & Observed & partial & MsgContext.Body = 'How to make a cake?'; MsgContext.SessionKey = agent:main:dashboard:trace-...; MsgContext.AgentId = main; MsgContext.ChatType = direct; MsgContext.MessageSid = 8d07c88e-188f-4347-9f9a-7fc90... \\
G14 & entered & Observed & partial & DispatchInboundResult = returned \\
G15 & finalized & Observed & partial & FinalizedMsgContext.Body = 'How to make a cake?'; FinalizedMsgContext.BodyForAgent = 'How to make a cake?'; FinalizedMsgContext.BodyForCommands = 'How to make a cake?'; FinalizedMsgContext.SessionKey = agent... \\
G16 & normal\_path & Observed & partial & DispatchFromConfigResult = returned \\
G17 & resolved & Observed & partial & sessionAgentId = main; sessionAgentCfg = not directly logged \\
G18 & selected & Observed & partial & replyResult = returned to G16 \\
\bottomrule
\end{longtable}


\subsection{G0--G3 Authentication and Authorization}
\label{app:r1-g0-g3}

The first four stages are the main difference from the old Cake2 trace.
For R1:

\begin{itemize}
    \item G0 records token authentication state with
    \texttt{sharedAuthProvided=True}, \texttt{sharedAuthOk=True}, and no
    device identity.
    \item G1 records the shared credential outcome as \texttt{allow}.
    \item G2 records role \texttt{operator} and scope
    \texttt{operator.write}; no fallback credential path is used.
    \item G3 records method authorization for \texttt{chat.send} with
    \texttt{operator.write}. The administrator shortcut is not an R1
    claim.
\end{itemize}

The source-level rule is narrower than a causal claim about
authentication method: \texttt{operator.admin} can short-circuit the
method-scope check; otherwise \texttt{chat.send} requires
\texttt{operator.write}. The controlled R1 trace and the historical old
Cake2 trace differ in observed scopes, but these cases do not prove that
authentication method alone determines scope.

\section{Full Stage-by-Run Matrix}
\label{app:stage-matrix}

\begin{table}[H]
\centering
\caption{Full compressed G0--G18 matrix for R1--R7. All controlled runs
share the same branch labels; each stage is directly observed at outcome
level while rationale visibility remains partial.}
\label{tab:app-stage-matrix}
{\scriptsize
\begin{tabularx}{\linewidth}{@{}l l l X@{}}
\toprule
Stage & Branch & Runs & Rationale visibility \\
\midrule
G0 & authorized & all R1--R7 & partial \\
G1 & allow & all R1--R7 & partial \\
G2 & pass & all R1--R7 & partial \\
G3 & allow & all R1--R7 & partial \\
G4 & pass & all R1--R7 & partial \\
G5 & unchanged & all R1--R7 & partial \\
G6 & resolved & all R1--R7 & partial \\
G7 & resolved & all R1--R7 & partial \\
G8 & pass & all R1--R7 & partial \\
G9 & resolved & all R1--R7 & partial \\
G10 & allow & all R1--R7 & partial \\
G11 & pass & all R1--R7 & partial \\
G12 & admitted & all R1--R7 & partial \\
G13 & constructed & all R1--R7 & partial \\
G14 & entered & all R1--R7 & partial \\
G15 & finalized & all R1--R7 & partial \\
G16 & normal\_path & all R1--R7 & partial \\
G17 & resolved & all R1--R7 & partial \\
G18 & selected & all R1--R7 & partial \\
\bottomrule
\end{tabularx}

}
\end{table}

\section{Agent Runtime Ledger and Metrics}
\label{app:ar-ledger}

The Agent Runtime layer is scored separately from the Gateway layer.
AR0--AR6 are seven logical stages:

\begin{longtable}{@{}p{0.10\linewidth}p{0.27\linewidth}p{0.55\linewidth}@{}}
\caption{AR-layer logical stages.}\label{tab:app-ar-definitions}\\
\toprule
Stage & Evidence type & Definition \\
\midrule
\endfirsthead
\toprule
Stage & Evidence type & Definition \\
\midrule
\endhead
AR0 & Runtime & Runtime attempt selected. \\
AR1 & Runtime & Agent lifecycle started. \\
AR2 & Runtime coverage & Either tool start/result pairs are observed, or a completed no-tool run records the no-tool outcome. \\
AR3 & Runtime & Final reply is observed at the Agent boundary. \\
AR4 & Runtime & Agent lifecycle ends. \\
AR5 & Runtime & Reply resolver returns a result. \\
AR6 & Source-derived & Later Gateway post-processing follows from the dispatch source after resolver return unless separately instrumented. \\
\bottomrule
\end{longtable}

\begin{table}[H]
\centering
\caption{AR-layer observations and tool-call-level ratios.}
\label{tab:app-ar-metrics}
{\scriptsize
\begin{tabularx}{\linewidth}{@{}l c l c l@{}}
\toprule
Run & AR observed & AR2 outcome & Tool result ratio & AR6 \\
\midrule
R1 & 6/7 & Observed(no\_tool) & n/a & SourceDerived \\
R2 & 6/7 & Observed(tool\_results) & 6/6 & SourceDerived \\
R3 & 6/7 & Observed(tool\_results) & 5/5 & SourceDerived \\
R4 & 6/7 & Observed(tool\_results) & 4/4 & SourceDerived \\
R5 & 6/7 & Observed(no\_tool) & n/a & SourceDerived \\
R6 & 6/7 & Observed(no\_tool) & n/a & SourceDerived \\
R7 & 6/7 & Observed(no\_tool) & n/a & SourceDerived \\
\bottomrule
\end{tabularx}

}
\end{table}

\section{Execution Classes}
\label{app:execution-classes}

\begin{table}[H]
\centering
\caption{Execution classes derived only from observed runtime paths.}
\label{tab:app-execution-classes}
{\scriptsize
\begin{tabularx}{\linewidth}{@{}l X l X@{}}
\toprule
Run & Prompt & Class & Observed tools \\
\midrule
R1 & How to make a cake? & no-tool & none \\
R2 & Search the web for the latest OpenTelemetry trac... & read-only retrieval & \texttt{web\_search} \( \times 6 \) \\
R3 & Draft an email to <redacted-email> saying that t... & external side effect via host app & \texttt{bash} \( \times 5 \) \\
R4 & Create a short Python script that prints the fir... & local file write & \texttt{bash} \( \times 3 \) + \texttt{apply\_patch} \\
R5 & Explain the difference between tracing and prove... & no-tool & none \\
R6 & Explain the difference between tracing and prove... & no-tool & none \\
R7 & How to make a cake? & no-tool & none \\
\bottomrule
\end{tabularx}

}
\end{table}

R2 is read-only retrieval: its runtime layer records six
\texttt{web\_search} calls with matching results. R4 is a local mutation:
it records shell activity, an \texttt{apply\_patch} call, and a follow-up
check inside the workspace. R3 is an external side-effect path: four
short capability probes are followed by one Apple Mail invocation through
\texttt{osascript}. The trace records exit code 0 for the host-side
command; downstream delivery is not verified. The submitted paper and
artifact must redact the recipient address and sender name.

\section{Exact-Prompt Repeatability}
\label{app:repeatability}

\begin{table}[H]
\centering
\caption{Exact-prompt repeatability for the two repeated prompt pairs.}
\label{tab:app-repeatability}
{\scriptsize
\begin{tabularx}{\linewidth}{@{}l l c c c c@{}}
\toprule
Pair & G0--G18 branch diff & Agent & Resolver & Tools & AR structure \\
\midrule
R1/R7 & none & yes & yes & yes & yes \\
R5/R6 & none & yes & yes & yes & yes \\
\bottomrule
\end{tabularx}

}
\end{table}

For R1/R7 and R5/R6, the repeated prompt pairs have no G0--G18 branch
differences. Agent, resolver, tool/no-tool choice, stop reason, and AR
structure also match. The accepted differences are run ID, Session
key/ID, timestamps, and reply wording.

\section{Historical Notes}
\label{app:historical}

\subsection{Old Cake2 Artifact}

The old Cake2 trace is retained in the artifact only as historical
context. It was collected under an earlier instrumentation version, took
a device-token/operator-admin path, and is excluded from all quantitative
results in this paper. It should not be used to compute R1--R7 coverage,
observation, repeatability, or AR metrics.

\subsection{Weather Runs}

The Weather runs are also historical notes. They showed that post-G18
tool behavior can vary within a prompt family: some runs recorded
\texttt{web\_search} only, while W3, W5, and the Beijing run also
recorded a \texttt{bash} skill-loading probe before \texttt{web\_search}.
They are useful artifact evidence for how post-G18 tool paths can differ
while G0--G18 remains stable, but they are not mixed with the controlled
R1--R7 scores.

\section{Privacy and Submission Boundary}
\label{app:privacy}

The working \texttt{artifacts/audit/} directory contains raw audit notes,
local paths, and sensitive prompt material. It is not a submission bundle.
Submitted artifacts must redact personal email addresses, sender names,
usernames, hostnames, absolute local paths, and author-identifying URLs.
If an audit-derived file is copied into a submission artifact, it must be
redacted first.
